\exercisesection{Classification}

\begin{enumerate}
\def\labelenumi{\arabic{enumi})}
\tightlist
\item
  \begin{enumerate}
  \def\labelenumii{\alph{enumii})}
  \tightlist
  \item
    For G to have a countable base, it is necessary and sufficient that \(\widehat{G}\) have a countable base. (If G has a countable base, then \(L^{1}(G)\) has a countable base, hence \(\mathsf{X}(\mathsf{L}^{1}(\mathsf{G}))\) has a countable base.)
  \end{enumerate}
\end{enumerate}

\begin{enumerate}
\def\labelenumi{\alph{enumi})}
\setcounter{enumi}{1}
\item
  For G to be countable at infinity, it is necessary and sufficient that \(\widehat{\mathrm{G}}\) admit an open subgroup with a countable base, or equivalently that \(\widehat{\mathrm{G}}\) be metrizable. (Use Prop. 3 of II, p.~248.)
\item
  Suppose G is compact. For G to be metrizable, it is necessary and sufficient that \(\widehat{G}\) be countable, or equivalently that G be isomorphic to a closed subgroup of \(T^{N}\) . (A countable discrete commutative group is a quotient of \(\mathbf{Z}^{(N)}\) .)
\end{enumerate}

\begin{enumerate}
\def\labelenumi{\arabic{enumi})}
\setcounter{enumi}{1}
\tightlist
\item
  \begin{enumerate}
  \def\labelenumii{\alph{enumii})}
  \tightlist
  \item
    Every infinite commutative group \(\Gamma\) admits an infinite countable quotient group. (Embed an infinite countable subgroup \(\Delta\) of \(\Gamma\) in a countable divisible group \(\Delta'\) by applying A, II, p.~185, Exercise 14; then extend to \(\Gamma\) the identity morphism of \(\Delta\) into \(\Delta'\) .)
  \end{enumerate}
\end{enumerate}

\begin{enumerate}
\def\labelenumi{\alph{enumi})}
\setcounter{enumi}{1}
\tightlist
\item
  Every infinite compact commutative group contains an infinite metrizable closed subgroup. (Use a).)
\end{enumerate}

\begin{enumerate}
\def\labelenumi{\arabic{enumi})}
\setcounter{enumi}{2}
\tightlist
\item
  \begin{enumerate}
  \def\labelenumii{\alph{enumii})}
  \tightlist
  \item
    Let K be a compact group and \(\mathrm{K}_{(n)}\) the set of elements of order n of K. If \(K \neq E_{n}\) for every n, then the set of elements of infinite order of K is dense in K. (If \(\mathrm{K}_{(n)}\) has nonempty interior, we have \(\mathrm{K}_{(m)} = \mathrm{K}\) for some integer \(m \geqslant n\) .)
  \end{enumerate}
\end{enumerate}

\begin{enumerate}
\def\labelenumi{\alph{enumi})}
\setcounter{enumi}{1}
\item
  Let \(\Gamma\) be a discrete commutative group whose elements are not of bounded order. Then \(\Gamma\) admits a countable quotient with the same property. (Reason as in part a) of Exercise 2.)
\item
  If every neighborhood of e in G contains an element of infinite order, then G possesses a metrizable closed subgroup with the same property. (Reduce to the case where G is compact, and use a) and b).) Otherwise, if G is not discrete, there exists an integer q \textgreater{} 1 such that G contains a closed subgroup isomorphic to \((\mathbf{Z}/q\mathbf{Z})^{\mathbf{N}}\) . (Use A, VII, p.~54, exerc. 4, c).)
\end{enumerate}

\begin{enumerate}
\def\labelenumi{\arabic{enumi})}
\setcounter{enumi}{3}
\item
  If G is compact and if \(\widehat{G}\) has an element of infinite order, then there exists a nontrivial continuous homomorphism from R to G. (Since R is divisible, there exists a nontrivial homomorphism from \(\widehat{G}\) into R.)
\item
  Let p be a prime number. The group \(Q_{p}\) is not the product of a compact group and a discrete group. (Every compact subgroup of \(Q_{p}\) is of the form \(p^{n}Z_{p}\) for some integer \(n \in Z\) .)
\item
  If G is a torsion group, then G and \(\widehat{G}\) are totally disconnected. (\(\widehat{G}\) is totally disconnected by Cor. 2 of II, p.~250; to prove that G is totally disconnected, reduce to the case where G is compact using Prop. 3 of II, p.~248, and then prove that \(\widehat{G}\) is a torsion group.)
\item
  An element \(x\) of \(G\) is said to be of infinite height if, for every \(n \in \mathbf{Z}\) , there exists \(y \in G\) such that \(y^n = x\) . Prove that, if \(G\) is compact, the set of elements of infinite height is the neutral component of \(G\) .
\end{enumerate}

¶ 8) The group G is locally connected if and only if G is isomorphic to a group \(\mathbf{R}^{n}\times\mathrm{E}\times\widehat{\mathrm{D}}\) where E and D are discrete and every finite-rank subgroup of D is free.

\begin{enumerate}
\def\labelenumi{\arabic{enumi})}
\setcounter{enumi}{8}
\tightlist
\item
  \begin{enumerate}
  \def\labelenumii{\alph{enumii})}
  \tightlist
  \item
    Let \(\boldsymbol{a}=(a_{n})_{n\geqslant0}\) be a sequence of integers \textgreater{} 1. We equip Z with the topological ring structure for which the subsets \(Za_{0}a_{1}\cdots a_{n}\) , where \(n\geqslant0\) , form a fundamental system of neighborhoods of 0. We denote by \(\Delta_{a}\) the completion of this ring. It is compact, metrizable, and totally disconnected. If p is a prime number, and if \(a_{i}=p\) for all i, we have \(\Delta_{a}=Z_{p}\) .
  \end{enumerate}
\end{enumerate}

\begin{enumerate}
\def\labelenumi{\alph{enumi})}
\setcounter{enumi}{1}
\item
  We denote by \(\mathbf{Z}(\boldsymbol{a}^{\infty})\) the multiplicative group of complex numbers of the form \(\exp(2i\pi l/(a_{0}a_{1}\cdots a_{r}))\) where \(l \in Z\) and \(r \in N\) , equipped with the discrete topology. Then \(\chi \mapsto \chi(1)\) is an isomorphism from \(\widehat{\Delta}_{\boldsymbol{a}}\) onto \(\mathbf{Z}(\boldsymbol{a}^{\infty})\) .
\item
  Let B be the subgroup of \(\mathbf{R} \times \Delta_{\boldsymbol{a}}\) consisting of the \((n, n)\) for \(n \in \mathbf{Z}\) ; it is discrete. We set \(\Sigma_{\boldsymbol{a}} = (\mathbf{R} \times \Delta_{\boldsymbol{a}})/\mathrm{B}\) . It is a compact, metrizable, connected group. (Note that the canonical image of \(\mathbf{R}\) into \(\Sigma_{\boldsymbol{a}}\) is dense, and use Lemma 1 of II, p.~244.)
\end{enumerate}

\(d)\) Let \(\Gamma_{\boldsymbol{a}}\) be the additive subgroup of the discrete group \(\mathbf{Q}\) consisting of rational numbers of the form \(l/(a_{0}a_{1}\cdots a_{r})\) where \(l\in\mathbf{Z}\) and \(r\in\mathbf{N}\) . If \(\chi\in\widehat{\Sigma}_{\boldsymbol{a}}\) , then \(\chi\) defines a unitary character of \(\mathbf{R}\times\Delta_{\boldsymbol{a}}\) , hence a unitary character of \(\mathbf{R}\) of the form \(x\mapsto\exp(2i\pi\alpha_{\chi}x)\) where \(\alpha_{\chi}\in\mathbf{R}\) . Then \(\chi\mapsto\alpha_{\chi}\) is an isomorphism from \(\widehat{\Sigma}_{\boldsymbol{a}}\) onto \(\Gamma_{\boldsymbol{a}}\) .

\begin{enumerate}
\def\labelenumi{\alph{enumi})}
\setcounter{enumi}{4}
\item
  Show that \(\widehat{\mathbf{Q}}\) is canonically isomorphic to \(\Sigma_{\pmb{a}}\) , where the sequence \(\pmb{a}\) is defined by \(a_{n} = n + 2\) for all \(n \geqslant 0\) .
\item
  Let \(R_{d}\) be the group R equipped with the discrete topology. Show that \(\widehat{R}_{d}\) is isomorphic to \((\widehat{\Sigma}_{\boldsymbol{a}})^{\mathfrak{c}}\) where \(a_{n}=n+2\) and where c is the cardinality of the continuum. (Consider a basis of R as a Q-vector space.)
\end{enumerate}

¶ 10) A topological group is said to be monothetic if there exists an element of the group whose powers are dense in the group, and solenoidal if there exists a continuous homomorphism from R into the group whose image is dense.

\begin{enumerate}
\def\labelenumi{\alph{enumi})}
\item
  With the notation of Exercise 9, the group \(\Delta_{\pmb{a}}\) is monothetic.
\item
  Every compact monothetic totally disconnected group is topologically isomorphic to a group of the form \(\Delta_{a}\) .
\item
  Suppose G is compact. For G to be monothetic, it is necessary and sufficient that \(\widehat{G}\) is isomorphic to a subgroup of T equipped with the discrete topology.
\item
  Suppose G is compact. The following conditions are equivalent:
\end{enumerate}

\begin{enumerate}
\def\labelenumi{(\roman{enumi})}
\item
  G is solenoidal;
\item
  \(\widehat{\mathbf{G}}\) is isomorphic to a subgroup of \(\mathbf{R}\) endowed with the discrete topology;
\item
  \(\widehat{\mathbf{G}}\) is torsion-free, and \(\operatorname{Card}(\widehat{\mathbf{G}}) \leqslant \mathfrak{c}\), the cardinality of the continuum;
\item
  G is a quotient of \((\Sigma_{\boldsymbol{a}})^{\mathfrak{c}}\) with \(\boldsymbol{a}=(n+2)_{n\geqslant0}\) .
\end{enumerate}

\begin{enumerate}
\def\labelenumi{\alph{enumi})}
\setcounter{enumi}{4}
\tightlist
\item
  If G is a compact solenoidal group, then G is monothetic.
\end{enumerate}

\begin{enumerate}
\def\labelenumi{\arabic{enumi})}
\setcounter{enumi}{10}
\tightlist
\item
  We denote by L(G) the set of continuous representations of G into R. A one-parameter subgroup of G is the image of a continuous homomorphism from R into G.
\end{enumerate}

\begin{enumerate}
\def\labelenumi{\alph{enumi})}
\tightlist
\item
  The following conditions are equivalent:
\end{enumerate}

\begin{enumerate}
\def\labelenumi{(\roman{enumi})}
\item
  G is the union of one-parameter subgroups;
\item
  every homomorphism from Z into G extends to a continuous homomorphism from R into G;
\item
  every continuous homomorphism from \(\widehat{\mathbf{G}}\) into \(\mathbf{R} / \mathbf{Z}\) arises by passing to the quotient from an element of \(\mathrm{L}(\widehat{\mathrm{G}})\) ;
\item
  every unitary character of \(\widehat{\mathbf{G}}\) is of the form \(e^{i\lambda}\) where \(\lambda \in \mathrm{L}(\widehat{\mathbf{G}})\) .
\end{enumerate}

\begin{enumerate}
\def\labelenumi{\alph{enumi})}
\setcounter{enumi}{1}
\item
  If G has a countable base, the conditions of a) are equivalent to: G is isomorphic to \(R^{n} \times T^{I}\) where I is a countable set.
\item
  The union of the one-parameter subgroups of G is a dense subgroup of the neutral component of G. (Use Exerc. 4.)
\end{enumerate}

\(d)\) Let \(t \mapsto \chi_t\) be a continuous map from \([0,1]\) into \(\widehat{\mathbf{G}}\) such that \(\chi_0\) is the trivial character. There exists a map \(t \mapsto \lambda_t\) from \([0,1]\) into \(\mathrm{L}(\mathrm{G})\) with \(\lambda_0 = 0\) such that, for all \(t\), we have \(\chi_t = \exp(2i\pi\lambda_t)\) , and, for every \(x \in \mathrm{G}\) , the map \(t \mapsto \lambda_t(x)\) is continuous.

\begin{enumerate}
\def\labelenumi{\alph{enumi})}
\setcounter{enumi}{4}
\tightlist
\item
  The union of the one-parameter subgroups of G is also the union of the arcs of G that contain e, an arc of G being the image of a continuous map from \([0,1]\) into G. (Use d).
\end{enumerate}

\begin{enumerate}
\def\labelenumi{\arabic{enumi})}
\setcounter{enumi}{11}
\tightlist
\item
  We use the notation L(G) from Exercise 11.
\end{enumerate}

\begin{enumerate}
\def\labelenumi{\alph{enumi})}
\item
  Let H be a closed subgroup of G. Every element of L(H) extends to an element of L(G). (First reduce to the case where \(G = R^{n} \times D\) with D discrete, then in the case where the given element of L(H) is trivial on \(R^{n}\) and then to the case where G is discrete. Then use the fact that R is divisible.)
\item
  Every continuous morphism from H into \(R^{n} \times T^{I}\) , where I is an arbitrary set, extends to a continuous morphism from G into \(R^{n} \times T^{I}\) (Use a) and Th. 4 of II, p.~226.)
\item
  Conversely, let A be a locally compact commutative group. Suppose that, for every locally compact commutative group G, for every closed subgroup H of G, and every continuous morphism \(\varphi\) of H into A, \(\varphi\) extends to a continuous morphism from G into A. Then there exist an integer n and a set I such that A is isomorphic to \(R^{n} \times T^{I}\) (Taking G = R, H = Z, one sees that A is connected, hence of the form \(R^{n} \times \widehat{D}\) with D discrete. Show that D is a projective Z-module, hence free.)
\end{enumerate}

\begin{enumerate}
\def\labelenumi{\arabic{enumi})}
\setcounter{enumi}{12}
\tightlist
\item
  Let C be a compact subset of G and \(\varepsilon > 0\) .
\end{enumerate}

\begin{enumerate}
\def\labelenumi{\alph{enumi})}
\item
  Let \(\alpha\) be the Haar measure of G. There exists a compact subset D of G such that \(\alpha(\mathrm{CD}) \leqslant (1 + \varepsilon)\alpha(\mathrm{D})\) . (The case where \(\mathbf{G} = \mathbf{R}^n\) being immediate, we are reduced to the case where G admits a compact open subgroup H, and we may assume C is saturated with respect to H; we are then reduced to the case where G is discrete; we may next assume G is finitely generated, and finally \(\mathbf{G} = \mathbf{Z}^p\) .)
\item
  There exists \(k \in \mathrm{L}^{1}(\widehat{\mathrm{G}})\) such that \(\mathcal{F}(k) \in \mathcal{K}(\mathrm{G})\) and \(\mathcal{F}(k)(x) = 1\) for \(x \in C\) , with \(\|k\|_{1} \leqslant 1 + \varepsilon\) . (Take k of the form \(\lambda fg\) , where \(\lambda\) is a constant and where \(f, g \in \mathrm{L}^{2}(\widehat{\mathrm{G}})\) have as Fourier transforms the characteristic functions of suitable sets, using a).)
\end{enumerate}

¶ 14) a) Suppose that every neighborhood of e in G contains an element of infinite order. There exists a compact metrizable perfect totally disconnected subset P of G which is a Kronecker set (Exercise 15 of II, p.~265), and a nonzero diffuse measure concentrated on P. (Reduce to the case where G is metrizable using Exercise 3, c). Then imitate the construction of TG, IX, p.~64, Lemma 3, and at the same time use part e) of Exercise 15 of II, p.~265.)

\begin{enumerate}
\def\labelenumi{\alph{enumi})}
\setcounter{enumi}{1}
\item
  If \(\mathrm{G} = (\mathbf{Z}/q\mathbf{Z})^{\mathbf{N}}\) , there exists a compact perfect totally disconnected subset P of G which is of type \(K_{q}\) , and a nonzero diffuse measure concentrated on P.
\item
  Assume G is non-discrete. There exists an element \(\tau\) of \(\mathcal{M}^{1}(G)\) such that \(\|\mathcal{F}(\tau)\|_{\infty}<\varrho(\tau)\) , where \(\varrho(\tau)\) is the spectral radius in \(\mathcal{M}^{1}(G)\) . There exists a non-invertible element \(\tau'\) of \(\mathcal{M}^{1}(G)\) such that \(\mathcal{F}(\tau')\geqslant1\) everywhere on \(\widehat{G}\) . There exists a non-Hermitian character of \(\mathcal{M}^{1}(G)\) . (Use a), b), exerc. 3, c), and part f) of exerc. 15 of II, p.~265.)
\end{enumerate}

\(d)\) Deduce from \(c\) ) that \(\widehat{\mathrm{G}}\) , which is open in \(\mathsf{X}(\mathcal{M}^1 (\mathrm{G}))\) , is not dense in \(\mathsf{X}(\mathcal{M}^1 (\mathrm{G}))\) .

